\emph{The lattice half.} $R_n/L_f$ is an $\F_\ell$-vector space (Lemma~\ref{lem:C}), so $[R_n:L_f]$ is a power of the odd prime $\ell$ and Lemma~S1 (the group lemma) applies with $k=2^{t-1}$: $L_f\cap2^{t-1}R_n=2^{t-1}L_f$; likewise $\ell R_n\cap2^{t-1}R_n=2^{t-1}\ell R_n$. Under (i), $L_f\cap J_t=L_f\cap2^{t-1}R_n=2^{t-1}L_f$, hence $\Lambda_{f,t}=\ell^{-1}H_n(2^{t-1}L_f)=2^{t-1}\Lambda_f$, and the vectors of $\Lambda_{f,t}$ coming from $a\notin\ell R_n$ are exactly the $2^{t-1}$-dilates of those of $\Lambda_f$ coming from $a\notin\ell R_n$ (the same lemma for $\ell R_n$).

\emph{What the depth-$t$ variant needs.} The depth-$t$ variant of Theorem~A is the argument of Prop.~F with $\Lambda_{f_0}$ replaced by $\Lambda_{f_0,t}$ and the floor $L_n$ of Lemma~A$^+$ replaced by the depth-$t$ floor $L_{n,t}$ of Lemma~\ref{lem:depthfloor}. We make explicit the unit to which the floor is applied, in both branches of Prop.~F. Suppose $\ell\mid k_n$, let $f_0$ be the saturation component of Lemma~\ref{lem:E}, and let $v\in\Lambda_{f_0,t}$ be a nonzero vector; write $v=\ell^{-1}H_na$ with $a\in L_{f_0}\cap J_t$, $a\ne0$. Since $a\in J_t$, $u_a\equiv1\pmod{2^t\mathcal O_n}$.
(Non-base branch, $a\notin\ell R_n$.) By Lemma~\ref{lem:E}, $r_a\in\RE$ and $v$ is the log vector of $r_a$, so $\hgt(r_a)=|v|$; $r_a\ne\pm1$ because $v\ne0$. The unit $r_a$ is a formal $\ell$-th root: $r_a^{\,\ell}=u_a\equiv1\pmod{2^t\mathcal O_n}$. Lemma~\ref{lem:oddtransfer} (odd-power depth transfer, $x=r_a$) gives $r_a\equiv1\pmod{2^t\mathcal O_n}$, and Lemma~\ref{lem:depthfloor} gives $\hgt(r_a)\ge L_{n,t}$.
(Base branch, $a=\ell b$, $b\in R_n\setminus\{0\}$.) By Lemma~\ref{lem:Eprime}, $v=H_nb$ is the log vector of $u_b\in A_n\subseteq\RE$ and $u_b\ne\pm1$; here $u_b^{\,\ell}=u_{\ell b}=u_a\equiv1\pmod{2^t\mathcal O_n}$, so Lemma~\ref{lem:oddtransfer} ($x=u_b$) gives $u_b\equiv1\pmod{2^t\mathcal O_n}$ and Lemma~\ref{lem:depthfloor} gives $\hgt(u_b)\ge L_{n,t}$.
In both branches the unit carrying $v$ has height $|v|\ge L_{n,t}$. Therefore the depth-$t$ variant reaches its contradiction exactly when Blichfeldt supplies a nonzero $v\in\Lambda_{f_0,t}$ with $|v|<L_{n,t}$; without Lemma~\ref{lem:oddtransfer} the passage from the depth of $u_a=r_a^{\,\ell}$ (the quantity filtered by $J_t$) to the depth of $r_a$ (the unit whose height is bounded) would be unjustified.

\emph{Comparison.} By the lattice half, a nonzero $v\in\Lambda_{f_0,t}$ with $|v|<L_{n,t}$ is $v=2^{t-1}v'$ with a nonzero $v'\in\Lambda_{f_0}$ and $|v'|<L_{n,t}/2^{t-1}$, whereas Theorem~A (Prop.~F) needs a nonzero $v'\in\Lambda_{f_0}$ with $|v'|<L_n$. Now $L_{n,t}/2^{t-1}<L_n$ iff $\operatorname{arcsinh}(2^t)<2^{t-1}\log(2+\sqrt5)$ iff $2^t+\sqrt{4^t+1}<(2+\sqrt5)^{2^{t-1}}$, and the last inequality holds for every $t\ge2$ (Lemma~S-analytic, \lbl{F}: \path{ineq_t_general}: $2^t+\sqrt{4^t+1}<2^{t+1}+1\le4^t\le4^{2^{t-1}}<(2+\sqrt5)^{2^{t-1}}$). Hence every vector that satisfies the depth-$t$ requirement satisfies the requirement of Theorem~A, and the requirement of the depth-$t$ variant is strictly stronger: the direct filtered-lattice sufficient condition at depth $t$ is stricter than that of Theorem~A, i.e.\ the resulting direct filtered-lattice criterion is weaker. This is the only sense of ``weaker'' used in this paper; it concerns the direct application of the same Blichfeldt argument to the deeper congruence sublattice $\Lambda_{f,t}$ and says nothing about other $2$-adic refinements. Nothing is asserted about the truth of hypothesis (i) beyond the certified cases of Cor.~S1.
